\begin{lemma}[Envelope slopes of a subsolution]
Let $F$ be a finite-dimensional real inner-product space, let $I\subseteq\R$ be
open, let $\Omega=\{(x,t):t\in I\}$, and let $f$ be a real function on space-time
which is continuous on $\Omega$ and has subquadratic spatial growth on compact
time intervals of $I$ in the sense of \noderef{SubquadraticGrowth},
and assume in addition that $f$ is a viscosity subsolution of $F[f]=0$ on $\Omega$ in the
sense of \noderef{IsLimitSubOn}. Then, with the slopes of \noderef{EnvSlopeP}
and \noderef{EnvSlopeM}:
\begin{enumerate}
\item for every $z=(x,t)$ and all $0<\ell\le\ell'$ such that $t+\sigma\in I$ for
all $\sigma\in[0,\ell']$,
\[
  \Senv_\ell^+f(z)\le\Senv_{\ell'}^+f(z);
\]
\item for every $z=(x,t)$, every $h>0$ with $t+\sigma\in I$ for all
$\sigma\in[-h,h]$, and every $\varepsilon>0$, there is $\delta>0$ with
\[
  \Senv_\ell^-f(z)\le\Senv_h^+f(z)+\varepsilon\qquad(0<\ell<\delta).
\]
\end{enumerate}
Equivalently, in the limit form: for every $z=(x,t)$ and every $h_0>0$ with
$t+\sigma\in I$ for all $\sigma\in[-h_0,h_0]$, item (1) makes
$h\mapsto\Senv_h^+f(z)$ nondecreasing on $(0,h_0]$, so that
$\lim_{h\downarrow0}\Senv_h^+f(z)$ exists in $[-\infty,\infty)$, and item (2)
gives
\begin{equation}\label{eq:subsolution-slope-order}
  \limsup_{\ell\downarrow0}\Senv_\ell^-f(z)
       \le\lim_{h\downarrow0}\Senv_h^+f(z),
\end{equation}
the right side being allowed the value $-\infty$.
\end{lemma}
\begin{proof}
\emph{(1).} Fix $z=(x,t)$ and $0<\ell\le\ell'$ with $t+\sigma\in I$ for all
$\sigma\in[0,\ell']$. By \noderef{QuadEnvPathConvex} with $T=\ell'$, $d=0$ and
$x_0=x$, the function $g(\sigma)=\Penv_\sigma f(x,t)$ is convex on
$[0,\ell']$, and $g(0)=f(z)$ because $\Penv_0f=f$ in \noderef{QuadEnvP}. By
\noderef{EnvSlopeP},
\[
  \Senv_\sigma^+f(z)=\frac{g(\sigma)-g(0)}{\sigma}\qquad(\sigma>0),
\]
which is the slope of the chord of $g$ from $0$ to $\sigma$. For a convex
function the chord slope from a fixed left endpoint is nondecreasing in the right
endpoint, so $\Senv_\ell^+f(z)\le\Senv_{\ell'}^+f(z)$.

\emph{(2).} Fix $z=(x,t)$, $h>0$ with $t+\sigma\in I$ for all $\sigma\in[-h,h]$,
and $\varepsilon>0$.

For $0<\ell<h$, the time arguments $t$ and $t-\ell$ lie in $I$, so by part (1)
of \noderef{QuadEnvAttained} the infimum defining $\Menv_\ell f(z)$ is attained
at some $y_\ell\in F$; put $\widetilde z_\ell=(y_\ell,t-\ell)$. Attainment and
the formula of \noderef{QuadEnvM} give
$\Menv_\ell f(z)=f(\widetilde z_\ell)+|y_\ell-x|^2/(2\ell)$, so
\[
  f(z)-\Menv_\ell f(z)
   =f(z)-f(\widetilde z_\ell)-\frac{|y_\ell-x|^2}{2\ell}.
\]
On the other hand $x$, at time $(t-\ell)+\ell=t$, is an admissible competitor in
the supremum defining $\Penv_\ell f(\widetilde z_\ell)$ of \noderef{QuadEnvP}.
That a competitor's value is at most the envelope is not automatic: the envelope
is a supremum, and the supremum of a set that is empty or unbounded above is a
junk value, so attainment at one competitor does not by itself license
domination at the others. The licence is part (1) of
\noderef{EnvelopeCompetitorBounds}, whose first clause states exactly that
$\Penv_\ell f$ dominates every competitor value; we apply it at the center
$\widetilde z_\ell$ and the length $\ell>0$, its hypotheses being the continuity
and subquadratic growth of $f$ already assumed together with the memberships
$t-\ell\in I$ and $(t-\ell)+\ell=t\in I$, both of which hold. This gives
\[
  \Penv_\ell f(\widetilde z_\ell)
    \ge f(z)-\frac{|x-y_\ell|^2}{2\ell}.
\]
Subtracting $f(\widetilde z_\ell)$ from both sides and dividing by $\ell>0$,
the two displays above give, by \noderef{EnvSlopeM} and \noderef{EnvSlopeP},
\begin{equation}\label{eq:eso-cross}
  \Senv_\ell^-f(z)\le\Senv_\ell^+f(\widetilde z_\ell).
\end{equation}

By part (4) of \noderef{QuadEnvAttained}, applied at $z$ with the tolerance
$\eta$, every optimizer $y_\ell$ for $\Menv_\ell f(z)$ satisfies
$|y_\ell-x|\le\eta$ once $\ell$ is small; since also $t-\ell\to t$, this gives
$\widetilde z_\ell\to z$ as $\ell\downarrow0$.

Fix $\ell$ with $0<\ell<h$, and apply part (1) at the center
$\widetilde z_\ell=(y_\ell,t-\ell)$ with lengths $\ell\le h$. Its hypothesis
requires $(t-\ell)+\sigma\in I$ for all $\sigma\in[0,h]$, which holds because
$(t-\ell)+\sigma=t+(\sigma-\ell)$ and $\sigma-\ell\in[-h,h]$. Therefore
\[
  \Senv_\ell^+f(\widetilde z_\ell)\le\Senv_h^+f(\widetilde z_\ell).
\]
Combining with \eqref{eq:eso-cross},
\[
  \Senv_\ell^-f(z)\le\Senv_h^+f(\widetilde z_\ell)
   =\frac{\Penv_hf(\widetilde z_\ell)-f(\widetilde z_\ell)}{h}
   \qquad(0<\ell<h).
\]
Now $\zeta\mapsto\Penv_hf(\zeta)$ is continuous at $z$: part (2) of
\noderef{QuadEnvAttained}, whose hypotheses $h>0$, $t\in I$ and $t+h\in I$ hold,
gives joint continuity of $(\ell,\zeta)\mapsto\Penv_\ell f(\zeta)$ at
$(h,z)$, and restricting to the slice $\ell=h$ gives continuity in the center.
Also
$f$ is continuous at $z$ because $z\in\Omega$ and $\Omega$ is open. Hence
$\zeta\mapsto\Senv_h^+f(\zeta)$ is continuous at $z$, so there is $\delta_1>0$
such that $|\zeta-z|<\delta_1$ implies
$\Senv_h^+f(\zeta)\le\Senv_h^+f(z)+\varepsilon$. Choose $\delta\in(0,h)$ so
small that $|\widetilde z_\ell-z|<\delta_1$ for all $0<\ell<\delta$, which is
possible by the previous paragraph. Then
\[
  \Senv_\ell^-f(z)\le\Senv_h^+f(z)+\varepsilon\qquad(0<\ell<\delta),
\]
as required.

Finally we derive \eqref{eq:subsolution-slope-order} from (1) and (2). Fix
$z=(x,t)$ and $h_0>0$ with $t+\sigma\in I$ for all $\sigma\in[-h_0,h_0]$.

For $0<h\le h'\le h_0$ the hypothesis of (1) at the center $z$ with the lengths
$h\le h'$ asks that $t+\sigma\in I$ for all $\sigma\in[0,h']$, which holds
because $[0,h']\subseteq[-h_0,h_0]$; so (1) gives
$\Senv_h^+f(z)\le\Senv_{h'}^+f(z)$, i.e. $h\mapsto\Senv_h^+f(z)$ is a
nondecreasing real-valued function on $(0,h_0]$. A nondecreasing real-valued
function on $(0,h_0]$ converges as $h\downarrow0$, in $\overline\R$, to
\[
  L=\inf_{0<h\le h_0}\Senv_h^+f(z)\in[-\infty,\infty),
\]
the value $-\infty$ occurring exactly when it is unbounded below, and the value
$+\infty$ being impossible because $L\le\Senv_{h_0}^+f(z)<+\infty$. Hence
$\lim_{h\downarrow0}\Senv_h^+f(z)$ exists in $[-\infty,\infty)$ and equals $L$.

Now fix $h\in(0,h_0]$ and $\varepsilon>0$. The hypothesis of (2) at $z$ with
this $h$ asks that $t+\sigma\in I$ for all $\sigma\in[-h,h]$, which holds
because $[-h,h]\subseteq[-h_0,h_0]$; so (2) supplies $\delta>0$ with
$\Senv_\ell^-f(z)\le\Senv_h^+f(z)+\varepsilon$ for all $0<\ell<\delta$.
Therefore $\limsup_{\ell\downarrow0}\Senv_\ell^-f(z)\le
\Senv_h^+f(z)+\varepsilon$. Letting $\varepsilon\downarrow0$ gives
$\limsup_{\ell\downarrow0}\Senv_\ell^-f(z)\le\Senv_h^+f(z)$ for every
$h\in(0,h_0]$, and taking the infimum over $h\in(0,h_0]$ gives
\[
  \limsup_{\ell\downarrow0}\Senv_\ell^-f(z)\le L
   =\lim_{h\downarrow0}\Senv_h^+f(z),
\]
which is \eqref{eq:subsolution-slope-order}.
\end{proof}
